\section{区域社会与边缘网络：三国竞争的地方条件}
\label{sec:three-kingdoms-regional-society-frontiers}

三国的疆域并不是三块内部同质的领地。魏在二二〇年定都洛阳，将尚书、禁军和北方仓储接回旧都；都城选择说明中枢试图把漕运、任官和宫廷警卫放进同一节点，却不能由此推出洛阳宫城复用的细节或城市人口。蜀汉在二二五年南征后重新取得南中部分铜铁、赋役、人力与道路，成都与南中地方集团的关系仍须不断协调。孙吴则在二三〇年遣舟师求夷洲、亶洲，远航中的疾病与补给困境显示，沿海动员能力并不自动变成对岛屿社会的稳定控制。

\subsection{东北不是魏、吴之间的一张空白地图}

公孙渊二三二年遣使通孙吴，孙权册拜其为燕王；次年公孙渊杀吴使、转而向魏请和。跨海使节和册封确实一度打开魏东北后方的风险，但孙吴不能以海路提供持续陆上援军，使盟约暴露出明显限度。公孙氏的意图、馈赠数额和辽东控制范围皆有争议，较可靠的结论是，辽东地方政权能够利用魏、吴间的距离和信息不对称，直到近距离军事压力改变了选择。

二三八年司马懿攻襄平，公孙渊败死，魏取得辽东的名义控制。征服后行政恢复不应从“平定”二字推成即时稳定；道路距离、战争暴力与地方官民的重新安置仍是后续问题。二四四年至二四五年毌丘俭进攻高句丽、攻破丸都城，也是同一限制的例子：魏军可以短期深入山地并击破都城，却没有由此建立长期直接统治。关于行军路线、山城位置和双方纪年，须将《三国志》、后出的朝鲜史书与丸都山城考古并列校读，而不能让任何一种叙事替代其他材料。

\subsection{陇右、淮南与江汉的不同战场}

二四七年姜维乘陇西羌胡起事出兵，魏将郭淮等应战。将羌胡行动直接写成蜀汉盟军，会掩盖部众与魏郡县之间本有的地方政治目标；蜀、魏的战线正是借由这些地方冲突彼此牵动。二五三年诸葛恪围合肥新城，吴军因疫病、攻城补给和魏守备撤退，远征的损失随即动摇其在建业的政治信誉。疾病存在并不允许用史书修辞量化死亡，仍足以说明淮南的水陆地形不能由一次胜利或一道诏令克服。

二六四年钟会在成都兵变，降附军民与魏军共同卷入短暂混乱；西晋因此须重新安置蜀地军政。吞并蜀汉并没有立即获得地方网络。二六九年羊祜在襄阳经营屯戍、招纳和对吴关系，二七二年西陵督步阐降晋、陆抗围城收复峡口，说明江汉、三峡与武昌以西的水陆枢纽仍可让一名边将或一处城戍影响全国战略。招纳人数、围城过程和将领个人作用不宜按传记夸写，地方联络的脆弱性却是两方持续经营边防的共同条件。

\subsection{统一是接收区域，而非取消区域}

二七九年西晋以益州水军沿江东下，荆、淮军并进；次年王濬舰队抵建业，孙皓出降。沿江攻势得以汇合，表明晋已把益州、荆州、淮南的不同资源和道路组织为同一计划；兵力、舰队规模和各军协同的细节仍有争议。晋所接收的是孙吴的州郡、官民与地方运行方式，而非可以立刻用北方尺度整齐丈量的江南。三国的终局因而也提醒我们：区域社会并非战事的被动背景，它们决定了何种运输、招募、结盟和安置可以真正完成。

\subsection{材料的边界}

《三国志》魏、蜀、吴诸纪传提供政权和战争年序，裴松之注保存公孙氏与蜀地兵变等异说，须同陈寿正文分层使用；《晋书》、朝鲜史书与《资治通鉴》分别补出晋军、东北和编年视角。洛阳城址、丸都山城考古及区域文书能够校正城市、山地与基层行政的某些问题，却不能直接给出全境人口或统一后的社会态度。本节只以这些不同亮度的材料追踪区域网络如何限制和支撑三国竞争。
